\documentclass[tikz,border=10pt]{standalone}
\usepackage{ctex}
\usepackage{tikz}
\usepackage{amsmath}
\usepackage{amssymb}
\usetikzlibrary{shapes.geometric, arrows.meta, positioning, calc, fit, backgrounds, shadows}

\begin{document}
\begin{tikzpicture}[
    node distance=1.3cm and 2.2cm,
    >=Stealth,
    font=\small,
    % 定义节点样式
    startstop/.style={
        rectangle, 
        rounded corners=10pt, 
        minimum width=3cm, 
        minimum height=1cm,
        text centered,
        draw=black,
        fill=gray!25,
        line width=1.2pt,
        drop shadow={shadow xshift=1pt, shadow yshift=-1pt, opacity=0.3}
    },
    process/.style={
        rectangle,
        rounded corners=6pt,
        minimum width=4cm,
        minimum height=1cm,
        text centered,
        draw=black!70,
        fill=white,
        line width=0.9pt,
        drop shadow={shadow xshift=0.5pt, shadow yshift=-0.5pt, opacity=0.2}
    },
    decision/.style={
        diamond,
        aspect=2.2,
        minimum width=2.8cm,
        minimum height=1.2cm,
        text centered,
        draw=orange!70!black,
        fill=orange!10,
        line width=0.9pt,
        drop shadow={shadow xshift=0.5pt, shadow yshift=-0.5pt, opacity=0.2}
    },
    arrow/.style={
        ->,
        >=Stealth,
        line width=1pt,
        color=black!70
    },
    looparrow/.style={
        ->,
        >=Stealth,
        line width=1pt,
        dashed,
        color=blue!60!black
    },
    % GNN模块样式（无阴影）
    gnnbox/.style={
        rectangle,
        rounded corners=8pt,
        minimum width=3.5cm,
        minimum height=0.9cm,
        text centered,
        draw=blue!70!black,
        fill=blue!15,
        line width=1pt
    },
    % DRL模块样式（无阴影）
    drlbox/.style={
        rectangle,
        rounded corners=8pt,
        minimum width=3.5cm,
        minimum height=0.9cm,
        text centered,
        draw=red!70!black,
        fill=red!15,
        line width=1pt
    },
    % 环境模块样式（无阴影）
    envbox/.style={
        rectangle,
        rounded corners=8pt,
        minimum width=3.5cm,
        minimum height=0.9cm,
        text centered,
        draw=green!70!black,
        fill=green!15,
        line width=1pt
    }
]

% ========== 输入层 ==========
\node[startstop] (input) at (0, 9.0) {输入：信道状态 $\mathbf{g}_{mk}^{(t)}$};

% ========== 特征构建 ==========
\node[process] (feature) at (0, 7.5) {节点/边特征构建\\$\mathbf{f}_k^{(t)}, \mathbf{e}_{ik}^{(t)}$};

% ========== GNN模块（蓝色区域）==========
\node[gnnbox] (gnn1) at (0, 5.8) {GNN消息传递（$L$层）};
\node[gnnbox] (gnn2) at (0, 4.6) {图嵌入输出 $\mathbf{z}_k^{(t)}$};

\begin{scope}[on background layer]
    \node[fit={(-2.5, 6.2) (2.5, 4.0)}, inner sep=8pt, fill=blue!8, rounded corners=12pt, draw=blue!50, line width=1.2pt] (gnnarea) {};
    \node[blue!70!black, font=\bfseries\large] at (0, 6.7) {GNN特征提取模块};
\end{scope}

% ========== DRL模块（红色区域）==========
\node[drlbox] (actor) at (6.2, 5.8) {Actor网络 $\pi_\theta$};
\node[drlbox] (action) at (6.2, 4.6) {动作采样 $\mathbf{a}_t = [\delta_k^{(t)}]$};

\begin{scope}[on background layer]
    \node[fit={(4.0, 6.2) (8.4, 4.0)}, inner sep=8pt, fill=red!8, rounded corners=12pt, draw=red!50, line width=1.2pt] (drlarea) {};
    \node[red!70!black, font=\bfseries\large] at (6.2, 6.7) {DRL决策模块};
\end{scope}

% ========== 环境交互模块（绿色区域）==========
\node[envbox] (power) at (-6.2, 5.8) {功率控制 $\eta_k^{(t)} = \delta_k^{(t)} \cdot \eta_{\max}$};
\node[envbox] (reward) at (-6.2, 4.6) {奖励计算 $\mathcal{R}_t = \mathcal{R}_t^{\text{sum}} + \mathcal{R}_t^{\text{urllc}}$};

\begin{scope}[on background layer]
    \node[fit={(-8.4, 6.2) (-4.0, 4.0)}, inner sep=8pt, fill=green!8, rounded corners=12pt, draw=green!50, line width=1.2pt] (envarea) {};
    \node[green!70!black, font=\bfseries\large] at (-6.2, 6.7) {环境反馈模块};
\end{scope}

% ========== Critic评估 ==========
\node[drlbox] (critic) at (6.2, 3.2) {Critic网络 $V_\phi$};
\node[process] (advantage) at (6.2, 2.0) {优势函数估计\\$A_t = G_t - V_\phi(\mathbf{s}_t)$};

% ========== 循环判断 ==========
\node[decision] (loop) at (0, 3.2) {回合\\结束?};

% ========== 策略更新 ==========
\node[process] (ppo) at (0, 1.3) {PPO-Clip策略更新\\$\theta \leftarrow \theta + \alpha_\theta \nabla_\theta L^{\text{CLIP}}$};
\node[process] (value) at (0, 0.0) {值函数更新\\$\phi \leftarrow \phi - \alpha_\phi \nabla_\phi (V_\phi - G_t)^2$};

% ========== 输出 ==========
\node[startstop] (output) at (0, -1.3) {输出：收敛策略 $\pi_\theta^*$};

% ========== 连接线 ==========
\draw[arrow] (input) -- (feature);
\draw[arrow] (feature) -- (gnn1);
\draw[arrow] (gnn1) -- (gnn2);
\draw[arrow] (gnn2) -- (actor);
\draw[arrow] (actor) -- (action);
\draw[arrow] (action) -- (power);
\draw[arrow] (power) -- (reward);
\draw[arrow] (reward) -- (loop);
\draw[arrow] (gnn2) -- (critic);
\draw[arrow] (critic) -- (advantage);
\draw[arrow] (advantage) -| (ppo);
\draw[arrow] (loop) -- (ppo);
\draw[arrow] (ppo) -- (value);
\draw[arrow] (value) -- (output);

% 循环回到特征构建
\draw[looparrow] (loop.east) -- ++(2.0, 0) |- node[right, pos=0.25, font=\small] {否} (feature.east);

% 结束循环
\draw[arrow] (loop) -- node[left, font=\small] {是} (ppo);

% ========== 图例 ==========

\end{tikzpicture}
\end{document}
